\section{Maximum cell mass and inverse volumes at all scales}\label{lc:cell}

This section gives an independent proof of the uniform-sphere sample
lower bound in Corollary~\ref{lk:regression}. We track what the complete
boundary-state history reveals about the signal, subtracting a
correction built from local cell masses. The correction accounts for
concentration already present in the conditional signal law. Its
geometric input controls inverse simplex volumes without any
bounded-support assumption.

After the model and stopping reductions below,
Section~\ref{lc:cell-selection} establishes coercivity of the correction,
Section~\ref{lc:cell-inverse-section} proves the inverse-volume bound
without a support restriction, Section~\ref{lc:cell-fibers} compares
incoming and outgoing exact-observation fibers, and
Section~\ref{lc:cell-endgame} completes the information-growth and
decoding estimates.

The information identities below are the conditional relative-entropy chain rule; see \cite[Section~3.3]{Austin2020} for its standard Borel formulation. All fiber comparisons needed here are proved explicitly.

We use uniform probability \(\sigma\) on \(S^{d-1}\) for the signal,
written \(U\) when random. All logarithms are natural except in dyadic
scale definitions. Set \(L=\log(1/\epsilon)\), and suppose
\(M=M(d)=o(d^2)\) and \(0<\epsilon\le1/10\). Work in sufficiently
large dimension that \(M\le d^2\).

Assume uniform-prior angular success at least \(2/3\) in the model of
Section~\ref{l:model}. Fix the shared rule randomness with success
strictly greater than \(3/5\). Lemmas~\ref{l:prior-ac} and~\ref{l:borel}
replace those fixed rules by everywhere Borel rules for this prior.
Then fix the initial state and the independent transition and output
seeds with success at least \(3/5\). The resulting deterministic
program has the same state bound. Its output is determined by its
terminal state and index.

We may assume \(T<dL\). Otherwise the desired lower bound is immediate.
There are at most \((T+1)2^M\) possible output labels. By the cap bound
of Lemma~\ref{l:caps-endpoint}, the union of their successful signals
has mass at most \((T+1)2^M\epsilon^{d-1}\). Hence
\[
 \frac35\le (T+1)2^M\epsilon^{d-1}.
\]
Since \(L\ge\log10\) and \(T<dL\),
\(\log(T+1)\le\log(2d)+L\). Taking logarithms gives
\((d-2)L\le M\log2+\log(2d)+\log(5/3)\). Thus \(L=O(d)\) and
\(T=O(d^2)\) in the case under consideration. The \(3/5\) endpoint
in Lemma~\ref{l:caps-endpoint} also gives \(T>d/20\). These reductions
will bound the entropy of the complete boundary history; they do not
supply the precision-dependent block estimate.

\subsection{Selected cells and their entropy cost}\label{lc:cell-selection}

Use the parameters
\[
 n=k=\lfloor d/100\rfloor,\qquad \alpha=(d-1)/2.
\]
Group sample layers into blocks of size \(n\), completing the last
block with independent unused samples. In each block denote by \(A\)
the \(n\)-by-\(d\) matrix of sample rows and by \(b=AU\) the responses.
We keep, for the analysis, a history consisting of the messages at all
the block boundaries so far, starting with an empty history. Write
\(V\) for the incoming history and \(W=(V,Z)\) for the history after
the next message.

The message records the ending program state or, if the program halted
within the block, the layer of that block and final state; it can be a
dummy message in subsequent blocks. An initial state is fixed. This
scheme also permits recording if halted at the start. Thus conditional
on \(V\) there are at most \((n+2)2^M+1\) messages. Given \(V,A,b\)
the next message can be formed, and \(A\) is fresh and independent of
\(V,U\). The final history determines the output.

There are \(O(d)\) blocks by the preceding reduction. The message
alphabet therefore bounds the entropy of every history by \(O(d^3)\).
This bound will control the scale selectors below. The complete history
need not be available in program memory.

For a history \(V\), let \(p_v=\Pr(V=v)\), and write \(\mu_v\) for the conditional law of \(U\) given \(v\), with density \(f_v\) with respect to \(\sigma\). Throughout we ignore conditioning states of probability zero. These densities exist and can be taken bounded by \(1/p_v\), since the marginal law is \(\sigma\). Write \(I_V=I(U;V)=\mathbb E\log f_V(U)\); use the analogous notations for \(W\). We will track information with a correction for local concentration: a log cell-mass quantity \(J_V\) below will be subtracted from \(I_V\) in bounding its growth per block.

Define levels \(j\ge 0\) of cells of scale \(r_j=2^{1-j}\). Level 0 consists of the whole sphere. For \(j\ge 1\) use the sphere intersected with a partition into half-open grid cubes of side \(r_j/\sqrt d\). Denote the cell containing \(u\) by \(C_j(u)\); each cell has diameter at most \(r_j\) and \(\sigma\)-measure at most \((2r_j)^{d-1}\). For any of our probability laws \(\mu=\mu_v\) define
\[
 m_\mu(u)=\max_{j\ge 0} \frac{\mu(C_j(u))}{r_j^\alpha}.
\]
In fact \(m_\mu\ge 2^{-\alpha}\) and the terms for large \(j\) are bounded by \(2^{d-1} r_j^{d-1-\alpha}/p_v\). Let \(j_v(u)\) choose the first level giving the maximum. It is bounded for each \(v\) by
\[
 1+\frac{\log_2(1/p_v)+d-1+\alpha}{d-1-\alpha}.
\]
Together the selector and the selected cell define a partition for each \(v\), with parts \(D\) consisting of points having the same selector and selected cell. When speaking of the random part \(D\) we include its identification by this choice (given \(v\)). Denote by \(p_D=\mu_v(D)\) its probability, by \(j=j_v(U)\) the chosen level, and by \(p_C=\mu_v(C_j(U))\) the original cell probability. Then
\begin{equation}\label{lc:cell-eq1}
\mathbb E\log(p_C/p_D)\le H(j\mid V).
\end{equation}
Indeed conditional on \(v,j\), \(p_D\) is the selector probability given \(v\) times the probability of the cell given \(v,j\). The expected log of its ratio with \(p_C\), in the order in \eqref{lc:cell-eq1}, equals the surprisal of the selector given \(v\) minus a nonnegative relative entropy for the cell distributions (at level \(j\)).

We record some finiteness and entropy bounds. By the bound on the selector and \(\mathbb E\log(1/p_V)=H(V)\), Jensen's inequality gives \(H(j\mid V)=O(\log d)\). Also \(H(D\mid V)<\infty\); the log of the number of cells at level \(j\) meeting the sphere is \(O(d(1+j))\). For this count, the disjoint cubes are within the ambient ball of radius \(1+r_j\); the ambient unit ball has volume \(\le (C/\sqrt d)^d\), as seen for example by integrating \(e^{-\|x\|^2/2}\) on the ball of radius \(\sqrt d\). The same entropy bounds apply with \(W\) in place of \(V\). In addition, the bin
\[
 t=\lfloor\log_2 m_{\mu_W}(U)\rfloor
\]
has \(H(t\mid W)=O(\log d)\): the range for each \(w\) has size \(O(d+\log(1/p_w))\), by the selector bound, the lower bound \(2^{-\alpha}\) and \(\mu_w(C_j)\le 1\). The use of \(O(\log d)\) here and above uses \(H(V),H(W)\le O(d^3)\).

Set \(J_V=\mathbb E\log m_{\mu_V}(U)\). All such expectations are
finite by the preceding bounds. Our objective on each block is
\[
 (I_W-J_W)-(I_V-J_V)\le Cd.
\]
We first show that the correction removes at most half the information,
up to \(O(d)\). Once the block estimate is proved, this coercivity
will turn its iteration into a bound on the information in the final
history.

\begin{lemma}[Coercivity of the cell correction]\label{lc:cell-coercivity}
For the conditional laws and selected-cell partition just defined,
\begin{equation}\label{lc:cell-eq2}
J_V \le \frac{\alpha}{d-1} I_V + O(d).
\end{equation}
\end{lemma}
\begin{proof}
For the selected cell \(C\),
\[
 \log m_{\mu_v}(U)
 \le \frac{\alpha}{d-1}\log\frac{p_C}{\sigma(C)}+\alpha\log 2,
\]
using \(\sigma(C)\le (2r_j)^{d-1}\) and \(\log p_C\le 0\). The log ratio here is at most \(\log(p_D/\sigma(D))+\log(p_C/p_D)\). The expected first term is at most \(I_V\) by relative entropy under the partition for each \(v\); \eqref{lc:cell-eq1} bounds the other term. Since \(\alpha/(d-1)=1/2\), the correction removes at most half the information, up to an additive \(O(d)\).
\end{proof}

We will group points according to bounds on \(m_\mu\) when taking projections. For any restriction of \(\mu\) to a set with \(m_\mu\le m_0\), the resulting subprobability \(\eta\) satisfies
\begin{equation}\label{lc:cell-eq3}
\eta(B(x,r))\le C^d m_0 r^\alpha\qquad (x\in\mathbb R^d,\ r>0),
\end{equation}
where here \(B(x,r)\) is an ambient ball. For \(r\le 1\), take a level \(j\ge 1\) with \(r\le r_j\le 2r\); a ball of radius \(r\) meets at most \(C^d\) cubes, by their side length and the volume bound applied to radius \(r+r_j\). Each cube contributing mass has \(\mu\)-mass at most \(m_0 r_j^\alpha\). For \(r>1\) use the mass bound 1 and \(m_0\ge 2^{-\alpha}\) if the set is nonempty.

\subsection{Inverse volumes at arbitrary length scales}\label{lc:cell-inverse-section}

\begin{lemma}[Inverse volume without a support restriction]\label{lc:cell-inverse}
Let \(n=k=\lfloor d/100\rfloor\) and \(\alpha=(d-1)/2\), with \(d\) sufficiently large. Suppose a Borel probability \(\nu\) on \(\mathbb R^d\) satisfies \(\nu(B(x,r))\le r^\alpha\) for all \(x,r\). Fix any \(u\), draw \(z_1,\ldots,z_k\) independently from \(\nu\), and denote by \(S\) the \(k\)-volume of the parallelepiped of \(z_1-u,\ldots,z_k-u\) (the square root of the Gram determinant). There is an absolute constant \(C\) such that
\begin{equation}\label{lc:cell-eq4}
\mathbb E S^{-n}\le \exp(Cdk).
\end{equation}
\end{lemma}
\begin{proof}
The volume is positive almost surely, since \(\nu\) gives zero mass to affine spaces of dimension \(i<\alpha\). This last fact follows by covering any bounded portion by \(O(r^{-i})\) balls of radius \(r\) as \(r\downarrow 0\), and suffices by conditioning successively on the previous draws.

We order the vectors \(z_i-u\) greedily, taking at each step one of
largest distance to the span of those preceding. Every tuple admits
such an ordering. By exchangeability, it suffices to bound the
expectation restricted to the original order being greedy and multiply
by \(k!\); this also allows ties. In this order let \(\delta_i\) be
the distance to the preceding span, with the norm used for the first
vector. Then \(S=\prod_i\delta_i\), and the distances are nonincreasing.
Fix their length bins
\[
 2^{-\ell_i-1}<\delta_i\le 2^{-\ell_i},\qquad
 \ell_1\le\cdots\le\ell_k,\quad \ell_i\in\mathbb Z.
\]
We will first bound the probability of these bins and the greedy
conditions, then sum their contributions to the inverse volume.

Given the first \(i-1\) draws, let \(e_j\), \(j<i\), be their
translated vectors' successive Gram--Schmidt axes. Group these axes
according to \(h=\ell_j\). Greediness forces the projection of
\(z_i-u\) onto each group to have norm at most \(2^{-h}\): before
the first axis in that group was chosen, the residual norm of
\(z_i-u\) was no greater than the distance chosen at that step.
Its residual norm after all preceding axes is at most \(2^{-\ell_i}\).

In a group of dimension \(a_h\), cover the ball of radius \(2^{-h}\)
by a net of radius \(2^{-\ell_i}2^{-(\ell_i-h)}\). Euclidean ball
packing bounds its cardinality by
\((3\cdot 2^{2(\ell_i-h)})^{a_h}\). The distinct group indices are
integers \(h\le\ell_i\), so their squared net errors sum to at most
\[
 2^{-2\ell_i}\sum_{h\le\ell_i}2^{-2(\ell_i-h)}
 =\tfrac43\,2^{-2\ell_i}.
\]
Including the orthogonal residual, all allowed points \(z_i\) therefore
lie in at most
\(3^{i-1}2^{2\sum_{j<i}(\ell_i-\ell_j)}\) ambient balls of radius
\(2\cdot2^{-\ell_i}\). These necessary constraints involve only the
preceding draws and the fixed bins. The local ball bound gives their
conditional probability at most
\[
 3^{i-1}2^\alpha\, 2^{2\sum_{j<i}(\ell_i-\ell_j)-\alpha\ell_i}.
\]

We may multiply these bounds sequentially, imposing only the necessary
constraints just described and never conditioning preceding draws on
future greedy choices. Use the local estimate when \(\ell_i\ge0\)
and probability one otherwise. Put
\[
 P=\sum_i\max(\ell_i,0),\qquad N=\sum_i\max(-\ell_i,0).
\]
The double sum in the product exponent satisfies
\[
 \sum_{i:\ell_i\ge0}\sum_{j<i}(\ell_i-\ell_j)
 \le \sum_{i:\ell_i\ge0}\left(k\ell_i+
                         \sum_{j:\ell_j<0}(-\ell_j)\right)
 \le k(P+N).
\]
Indeed, each inner sum is \((i-1)\ell_i-\sum_{j<i}\ell_j\).
Its first term is at most \(k\ell_i\), subtracting nonnegative
\(\ell_j\) can only decrease it, and each negative index contributes
at most \(k\) times. The probability of the greedy and length conditions
is consequently at most
\[
 \exp(Cdk)\,2^{-(\alpha-2k)P+2kN}.
\]

On these bins, \(S^{-n}\le2^{nk}2^{n(P-N)}\). When negative bins
dominate, the large distances already make this inverse volume small;
otherwise the local probability bound supplies the needed decay.
For \(N\le2P\), combining the two bounds gives the exponent
\[
 -(\alpha-2k-n)P+(2k-n)N
 \le -(\alpha-6k-n)P-nN\le -(P+N),
\]
since \(\alpha-6k-n\ge3\) and \(n\ge3\) for large \(d\).
For \(N>2P\), use probability at most one: then
\(n(P-N)\le -(P+N)\), because
\((n+1)P\le(n-1)N\) for \(n\ge3\).
Thus either region contributes at most \(\exp(Cdk)2^{-(P+N)}\).
The sum of \(2^{-(P+N)}\), even over all integer tuples, is \(3^k\).
This factor and \(k!\) are absorbed by \(\exp(Cdk)\), proving
\eqref{lc:cell-eq4} at arbitrary length scales.
\end{proof}

The inverse-volume estimate now controls a projection density at a prescribed target, even when that target lies outside the support of the projected measure.

\begin{lemma}[A fixed-target Gaussian moment]\label{lc:cell-moment}
For a subprobability \(\eta\ll\sigma\) on the sphere satisfying \eqref{lc:cell-eq3}, let \(q=\eta(S^{d-1})>0\). For Gaussian \(A\) write \(g_{\eta,A}\) for the Lebesgue density of \(Az\) under \(\eta\). Such densities exist almost surely since \(A\) has rank \(n<d-1\) and they exist for \(\sigma\). We may evaluate a version by the liminf of mass densities on shrinking cubes (say radii \(1/l,\ l\to\infty\)), which gives the Lebesgue density almost everywhere by differentiation. With this latter evaluation, for every fixed \(u\) we have
\begin{equation}\label{lc:cell-eq5}
\mathbb E_A g_{\eta,A}(Au)^k
 \le \left[\exp(Cd) q^{1-n/\alpha}m_0^{n/\alpha}\right]^k.
\end{equation}
\end{lemma}
\begin{proof}
Take \(k\) independent draws from \(\eta/q\). By \eqref{lc:cell-eq3}, scaling the space (also scaling \(u\)) by \((C^d m_0/q)^{1/\alpha}\) makes \eqref{lc:cell-eq4} applicable. In the original space the bound for \(\mathbb E S^{-n}\) is multiplied by \((C^d m_0/q)^{kn/\alpha}\). For fixed draws, the probability that all \(\|A(z_i-u)\|_\infty\le a\) is bounded by
\((2a)^{nk}(2\pi)^{-nk/2} S^{-n}\), by the maximum Gaussian density in \(k\) coordinates for each row. The \(A\)-expectation of the cube mass to the power \(k\) is \(q^k\) times the draws' average of this event probability. Thus we get \eqref{lc:cell-eq5} for the mass divided by \((2a)^n\) in place of the density. Fatou's lemma proves it for the liminf. Constants use that \(n/\alpha\) is bounded.
\end{proof}

\subsection{Comparing incoming and outgoing fibers}\label{lc:cell-fibers}

We return to \(V,W,A,b\) on the block and write \(\beta=n/\alpha\).
We use the same incoming selected part in both conditional signal laws.
Its change in probability can then be charged to the information
gained by the block, as we spell out below. Thus all parts \(D\) in
this comparison use the partition from \(V\); let \(r=r_j\) be the
selected scale. Use these two densities evaluated at \(b\):
\begin{itemize}
\item \(g_{v,D,A}\) for the projection by \(A\) of \(\mu_v\) restricted to \(D\);
\item \(g_{w,D,t,A}\) for the projection of \(\mu_w\) restricted to \(D\cap\{z:\lfloor\log_2 m_{\mu_w}(z)\rfloor=t\}\).

\end{itemize}
As random quantities the indices here take their actual values. Put
\(q_D=\mu_w(D)\), and denote by \(q_{Dt}\) this last restriction's
mass. They are positive almost surely for the actual indices.
Because \(D\) is determined by \((V,U)\) and \(V\) is retained in
\(W=(V,Z)\), its probability change satisfies
\begin{equation}\label{lc:cell-mass-change}
 \mathbb E\log(q_D/p_D)=I(D;Z\mid V)
 \le I(U;Z\mid V)=I_W-I_V.
\end{equation}
The projection bounds will attach the coefficient \(1-\beta\) to
this term, leaving a positive fraction of the information gain to
control.

At fixed full rank \(A\), let \(g_{0,A}\) be the projected density
from \(\sigma\), and let \(\lambda_{A,b}\) be the corresponding
conditional probability on the fiber at \(b\). In orthonormal
row-space coordinates the \(\sigma\)-projection has density
\(c_{d,n}(1-\|x\|^2)^{(d-n-2)/2}\) in the unit ball, and the remaining
component is conditionally uniform on its sphere. These facts follow,
for example, by normalizing a standard Gaussian to generate \(\sigma\)
and splitting it into the two orthogonal spaces; the squared lengths
before normalization are independent chi-squares. Transforming the
row-space coordinates to \(b\) multiplies the density by
\(\det(AA^\top)^{-1/2}\). The fiber probabilities can be taken
measurably by adding to the fixed row-space component a suitably scaled,
normalized Gaussian projection in the null space.

For a restricted density \(f{\bf 1}_E\) relative to \(\sigma\), a
projected density version is
\[
 g_{0,A}(b)\int f{\bf 1}_E\,d\lambda_{A,b}
\]
where \(g_{0,A}>0\), and zero outside. We use this version for
conditional probability formulas. It agrees with the differentiation
version almost everywhere for Lebesgue \(b\) at almost every \(A\).
The versions also agree when needed in the actual variables: the law
of \((A,b)\) is absolutely continuous with respect to Gaussian \(A\)
times Lebesgue measure, and there are finitely many history choices
and countably many sets used here.

We first lower bound the expected log of the incoming restricted
projection density:
\begin{equation}\label{lc:cell-eq6}
\mathbb E\log g_{v,D,A}(b)
 \ \ge\ (1-\beta)\mathbb E\log p_D +\beta J_V
       -\beta H(j\mid V)-O(n).
\end{equation}
Given \(V,D,A\), the density of \(b\) is \(g_{v,D,A}/p_D\), since
\(A\) is independent of \(U,V\). Take \(c\) to be the grid cube center
(0 at level 0), so \(\|U-c\|\le r\). Compare this conditional density
with a Gaussian of center \(Ac\) and covariance \(r^2I_n\).
Nonnegativity of relative entropy gives
\[
 \mathbb E\log g_{v,D,A}(b)\ \ge\ \mathbb E[\log p_D-n\log r]-O(n).
\]
Here \(\mathbb E[\|A(U-c)\|^2/r^2\mid V,D]\le n\). The identity
\[
 \log p_D-n\log r
 =(1-\beta)\log p_D+\beta\log m_{\mu_v}(U)
                  -\beta\log(p_C/p_D)
\]
and the selector entropy bound \eqref{lc:cell-eq1} give
\eqref{lc:cell-eq6}.

\begin{lemma}[Comparison on the exact observation fiber]\label{lc:cell-fiber-comparison}
With the incoming selected part \(D\), outgoing mass bin \(t\), and density versions specified above,
\begin{equation}\label{lc:cell-eq7}
I_W-I_V
 \ \le\ \mathbb E\log g_{w,D,t,A}(b) - \mathbb E\log g_{v,D,A}(b)
       + H(t\mid W,D).
\end{equation}
\end{lemma}
\begin{proof}
Conditional on \(V,D,A,b,W\), the density of \(U\) relative to \(\lambda_{A,b}\) is
\[
 f_v{\bf 1}_D\, g_{0,A}(b)/g_{v,D,A}(b),
\]
since the message given \(V,A,b\) no longer depends on \(U\). Compare on this fiber with the subprobability density
\[
 \sum_{t'} \frac{q_{Dt'}}{q_D}\,
   f_w{\bf 1}_{D,\ t'}\, g_{0,A}(b)/g_{w,D,t',A}(b)
\]
where zero-denominator terms are omitted, and the indicator restricts
to the specified bin of \(m_{\mu_w}\) as well as to \(D\).
Each included fiber density before its mixing weight has integral
one. Omitting bins can only reduce the total mass, so the displayed
mixture is a subprobability.

On the actual draw, the term of its bin has positive denominator
almost surely. Indeed, a zero denominator with \(g_{0,A}(b)>0\)
makes \(f_w{\bf 1}_{D,\ t'}\) zero \(\lambda_{A,b}\)-almost everywhere.
The actual conditional law is absolutely continuous on the fiber,
and \(f_w(U)>0\) almost surely. The log ratio of the true density
to the comparison density is therefore
\[
 \log\frac{f_v(U)}{f_w(U)}
 +\log\frac{g_{w,D,t,A}(b)}{g_{v,D,A}(b)}
 -\log\frac{q_{Dt}}{q_D}.
\]
The last term averages to \(H(t\mid W,D)\). Once integrability is
verified, the nonnegative expectation of this log ratio gives
\eqref{lc:cell-eq7}, by Jensen on each fiber against the subprobability.

We verify the logarithmic integrability needed here and in
\eqref{lc:cell-eq6}. The log densities
\(\log f_V(U)\) and \(\log f_W(U)\) are integrable by the state probability
bounds and \(x\log^-x\le1/e\) for densities relative to a probability.
Both projection densities have integrable positive log parts: they
are bounded above by \(g_{0,A}/p_v\) or \(g_{0,A}/p_w\), and
\[
 g_{0,A}\le c_{d,n}\det(AA^\top)^{-1/2}
\]
has integrable positive log part. For this last assertion,
Gram--Schmidt on Gaussian rows gives squared successive distances
with conditional chi-square laws and integrable logarithms.

For the incoming density \(g_{v,D,A}\), the Gaussian comparison also
bounds the negative log part. A log density ratio under its numerator
law has expected negative part at most \(1/e\) when compared to a
probability. The Gaussian comparison density has integrable negative
log part by \(r\le2\) and the squared norm estimate, while
\(-\log p_D\) has finite expectation. Thus
\(\log g_{v,D,A}(b)\) is integrable.

The displayed fiber log ratio now has integrable positive part, using
the outgoing positive log bound, the incoming integrability, and the
finite expected surprisal for \(t\mid W,D\). Its expected negative
part is at most \(1/e\), which remains true for comparison with a
subprobability. The log ratio is therefore integrable, and its
displayed decomposition also gives integrability of
\(\log g_{w,D,t,A}(b)\). This justifies \eqref{lc:cell-eq7}.
\end{proof}

We have compared the two fiber laws. The remaining step is to charge the correlation between the outgoing history and this fresh Gaussian block; the earlier blocks will carry no additional charge here.

To upper bound the target projection term, change the joint law \(Q\) of \(U,W,A\) to a reference \(R\) that keeps the joint \(U,W\) but makes \(A\) independent Gaussian. The relative entropy cost satisfies
\[
 D_{\rm KL}(Q\|R)=I(A;U,W)=I(A;Z\mid U,V)\le M\log 2+O(\log d),
\]
by freshness and the message alphabet. Under \(R\), condition on \(U,W\) to fix all indices. Apply \eqref{lc:cell-eq5} to the indicated restriction of \(\mu_w\), using \eqref{lc:cell-eq3} with \(m_0=2^{t+1}\le 2 m_{\mu_w}(U)\) and mass \(q_{Dt}\le q_D\). Thus the variable
\[
 X=\frac{g_{w,D,t,A}(AU)}
         {q_D^{1-\beta}m_{\mu_w}(U)^\beta}
\]
has \(\mathbb E_R X^k\le \exp(Cdk)\). Here one can use the differentiation version in \eqref{lc:cell-eq5}, which suffices also by version agreement under \(Q\); \(\log X\) is integrable under \(Q\) by the bounds above (also \(H(D\mid W)\le H(D\mid V)\)). The entropy bound \(k\mathbb E_Q\log X\le D_{\rm KL}(Q\|R)+\log\mathbb E_R X^k\) follows from Jensen's inequality with the likelihood ratio. This is the locally justified entropy variational inequality \eqref{l:tilt}; its unbounded logarithm is integrable by the preceding fiber calculation. Since \(k\) is proportional to \(d\) and \(M\le d^2\), we have proved
\begin{equation}\label{lc:cell-eq8}
\mathbb E\log g_{w,D,t,A}(b)
 \ \le\ (1-\beta)\mathbb E\log q_D + \beta J_W + O(d).
\end{equation}

Combining \eqref{lc:cell-eq6}--\eqref{lc:cell-eq8} and the selector and bin entropy bounds gives
\[
 I_W-I_V\ \le\ (1-\beta)\mathbb E\log(q_D/p_D)
                    +\beta(J_W-J_V)+O(d).
\]
Since \(0<\beta<1\), the shared-part bound
\eqref{lc:cell-mass-change} can be absorbed on the left:
\[
 \beta(I_W-I_V)\le \beta(J_W-J_V)+O(d).
\]
The ratio \(\beta=n/\alpha\) tends to \(1/50\), so it is bounded
away from zero. Dividing by it proves the block estimate
\begin{equation}\label{lc:cell-eq9}
(I_W-J_W) - (I_V-J_V) \le O(d).
\end{equation}
Using incoming cells both in the incoming projection bound and in restricting the target is useful here: their probability change costs only a fraction of the information gain itself. Binning the target by local cell mass allows the random projection estimate without assuming an isotropic or single-scale conditional law.

\subsection{Information growth and decoding}\label{lc:cell-endgame}

Initially \(I_V=0\) and \(-J_V\le \alpha\log 2\). By \eqref{lc:cell-eq9}, and \eqref{lc:cell-eq2} with \(\alpha/(d-1)=1/2\), the final history has information at most
\[
 C d(1+\lceil T/n\rceil)\le C'(d+T).
\]
On the other hand compare its joint law with \(U\) to the product of their marginals on the event of chord error at most \(\epsilon\) for its output. The joint probability is at least \(3/5\), while the product probability is at most \(\epsilon^{d-1}\) by the cap bound. Information, by relative entropy on this binary event, is at least \((3/5)(d-1)L-\log 2\). It follows that \(T\ge c dL-C''d\) for absolute constants with \(c>0\). For large \(L\) this yields the desired bound directly, and for the remaining bounded \(L\) it follows by \(T>d/20\) shown earlier.

All block projections and fiber laws used here are for the exact, noiseless data. The cost constraint used on a block message follows from the bit-state bound even when computation on the data within transitions is unrestricted. This completes this independent information proof of the sample lower bound.
